% SGA 3, Expose II, Loop 1 English reconstruction.
% Scope: Section 3, immediately after Nota 2.2.2 through 3.14, stopping
% immediately before Section 4.
% Authority: Polo--Gille Exp2-14oct24.pdf, component-local pages 7--21,
% with the Section 4 boundary verified on p. 22 (combined-reader pp. 65--80).
% Force-OCR is a drafting/locator witness only.
% English comparison witness: Jacob C. Reinhold, jcreinhold/sga commit
% e7a259f3f8608ad3edf9bf6eead3fd504dd2d23e (CC BY 4.0); consulted for
% comparison and drafting, not treated as source authority.

% Source page: Exp2 component-local p. 7; combined-reader p. 65.
\section*{3. The tangent bundle, condition (E)}
\addcontentsline{toc}{section}{3. The tangent bundle, condition (E)}

Throughout this section, unless otherwise stated, the letters
\(\mathcal M,\mathcal M',\mathcal N,\ldots\) will denote finite free
\(\mathcal O_S\)-modules (that is, modules isomorphic to finite direct sums
of copies of \(\mathcal O_S\)).

We shall systematically use the identifications established in Expos\'e~I.
Thus, by a ``functor over \(S\)'' we shall mean indifferently either a
functor equipped with a morphism to \(S=h_S\), or a functor on the category
of objects over \(S\). We shall use ``group functor over \(S\),'' and
similar expressions, in the same way.

\subsection*{Definition 3.1.}
Let \(S\) be a scheme, let \(\mathcal M\) be a finite free
\(\mathcal O_S\)-module, and let \(X\) be a functor over \(S\). The
\emph{tangent bundle of \(X\) over \(S\), relative to the
\(\mathcal O_S\)-module \(\mathcal M\)}, is the \(S\)-functor
\[
  T_{X/S}(\mathcal M)
  =\underline{\operatorname{Hom}}_S(I_S(\mathcal M),X).
\]
In particular, the \emph{tangent bundle of \(X\) over \(S\)} is
\footnote{Editor's note: when \(S=\operatorname{Spec}(k)\) and \(X\) is a
\(k\)-scheme, one has
\(T_{X/S}(S')=\operatorname{Hom}_S(S'\otimes_k k[\epsilon],X)
=X(S'\otimes_k k[\epsilon])\); this recovers one of the usual definitions
of the tangent bundle.}
\[
  T_{X/S}=T_{X/S}(\mathcal O_S)
  =\underline{\operatorname{Hom}}_S(I_S,X).
\]

The assignment \(\mathcal M\mapsto T_{X/S}(\mathcal M)\) is a covariant
functor from finite free \(\mathcal O_S\)-modules to \(S\)-functors. In
particular, the morphisms \(\mathcal M\to0\) and \(0\to\mathcal M\)
define, respectively, an \(S\)-morphism
\[
  \pi_{\mathcal M}:T_{X/S}(\mathcal M)\longrightarrow
  T_{X/S}(0)\simeq X
\]
and a section \(\tau_0\) of it, called the \emph{zero section}.
\footnote{Editor's note: the original order ``\(0\to\mathcal M\) and
\(\mathcal M\to0\)'' has been corrected to ``\(\mathcal M\to0\) and
\(0\to\mathcal M\).''}

\subsection*{Remark 3.1.1.}
\footnote{Editor's note: paragraphs 3.1.1 and 3.1.2 have been added.}
The projection
\(\pi_{\mathcal M}:T_{X/S}(\mathcal M)\to X\) is induced by the zero
section \(\varepsilon_{\mathcal M}:S\to I_S(\mathcal M)\), whereas the
zero section \(\tau_0:X\to T_{X/S}(\mathcal M)\) is induced by the
structural morphism \(\rho:I_S(\mathcal M)\to S\). More explicitly, let
\(t\in T_{X/S}(\mathcal M)(S')\) correspond to the \(S\)-morphism
\(f:S'\times_S I_S(\mathcal M)\to X\), and let \(x\in X(S')\)
correspond to \(g:S'\to X\). Then
\[
  \pi(t)=f\circ(\operatorname{id}_{S'}\times\varepsilon_{\mathcal M}),
  \qquad
  \tau_0(x)=g\circ(\operatorname{id}_{S'}\times\rho).
\]

It follows from 3.1 that \(\mathcal M\mapsto T_{X/S}(\mathcal M)\) is a
covariant functor from finite free \(\mathcal O_S\)-modules to functors
over \(X\). In particular, \(\mathcal O(S)\) is an operator monoid of the
\(X\)-functor \(T_{X/S}(\mathcal M)\), compatible with functoriality in
\(\mathcal M\).

% Source page: Exp2 component-local p. 8; combined-reader p. 66.
\subsection*{Scholium 3.1.2.}
\footnotemark[15]
The preceding discussion means, in particular, the following. For every
\(S\)-morphism \(\varphi:X'\to X\), put
\[
  \Sigma(X',\mathcal M)
  =\operatorname{Hom}_X\bigl(X',T_{X/S}(\mathcal M)\bigr).
\]
The multiplicative monoid
\(\operatorname{End}_{\mathcal O_S}(\mathcal M)\) acts on
\(\Sigma(X',\mathcal M)\); write the action as
\((\lambda,x)\mapsto\lambda*x\). Then
\[
  \lambda*(\mu*x)=(\lambda\mu)*x,\qquad
  1*x=x,\qquad 0*x=\tau_0\circ\varphi.
\]
There is likewise an action of
\(\operatorname{End}_{\mathcal O_S}(\mathcal M\oplus\mathcal M)\) on
\(\Sigma(X',\mathcal M\oplus\mathcal M)\).
\footnote{Editor's note: in fact, only the actions of \(\mathcal O(S)\)
(respectively \(\mathcal O(S)\times\mathcal O(S)\)) on
\(\Sigma(X',\mathcal M)\) (respectively
\(\Sigma(X',\mathcal M\oplus\mathcal M)\)) will be used; see below and
the proof of 3.6.}

Let \(m:\mathcal M\oplus\mathcal M\to\mathcal M\) be addition, and let
\(\delta:\mathcal M\to\mathcal M\oplus\mathcal M\) be the diagonal.
Write
\[
  m_{X'}:\Sigma(X',\mathcal M\oplus\mathcal M)	o
    \Sigma(X',\mathcal M),
  \qquad
  \delta_{X'}:\Sigma(X',\mathcal M)	o
    \Sigma(X',\mathcal M\oplus\mathcal M)
\]
for the induced maps. For \(\lambda,\mu\in\mathcal O(S)\), let
\(\ell_\lambda\) denote multiplication by \(\lambda\) on \(\mathcal M\),
and let \(\ell_{\lambda,\mu}\) denote multiplication by
\((\lambda,\mu)\) on \(\mathcal M\oplus\mathcal M\). Since
\[
  m\circ\ell_{\lambda,\lambda}=\ell_\lambda\circ m,
  \qquad
  m\circ\ell_{\lambda,\mu}\circ\delta=\ell_{\lambda+\mu},
\]
we have, for every \(z\in\Sigma(X',\mathcal M\oplus\mathcal M)\) and
\(x\in\Sigma(X',\mathcal M)\),
\begin{equation}
  \lambda*m(z)=m((\lambda,\lambda)*z),
  \qquad
  m((\lambda,\mu)*\delta(x))=(\lambda+\mu)*x.
  \tag{$\dagger$}
\end{equation}

\subsection*{Definition 3.2.}
Let \(u\in X(S)=\operatorname{Hom}_S(S,X)=\Gamma(X/S)\). The
\emph{tangent space of \(X\) over \(S\) at \(u\), relative to
\(\mathcal M\)}, denoted \(L^u_{X/S}(\mathcal M)\), is the \(S\)-functor
obtained from the \(X\)-functor \(T_{X/S}(\mathcal M)\) by pullback along
\(u:S\to X\).

\SGAThreeDiagram[0.36\textwidth]
  {figures/exp2/exp2_p008_tangent_space_cartesian_square.png}
  {The cartesian square defining \(L^u_{X/S}(\mathcal M)\).}
  {fig:exp2-tangent-space-cartesian}
  {Polo--Gille Exp2-14oct24.pdf, component-local p. 8; combined-reader p. 66}

In particular, \(L^u_{X/S}(\mathcal O_S)\) is denoted \(L^u_{X/S}\) and
is called the \emph{tangent space of \(X\) over \(S\) at \(u\)}.

\subsection*{Remark 3.2.1.}
\footnote{Editor's note: this remark has been added.}
It follows from 3.1.1 that, for every \(t:S'\to S\),
\(L^u_{X/S}(\mathcal M)(S')\) is the set of \(S\)-morphisms
\(f:I_{S'}(\mathcal M)\to X\) such that
\(f\circ\varepsilon_{\mathcal M}=u\circ t\), where
\(\varepsilon_{\mathcal M}:S'\to I_{S'}(\mathcal M)\) is the zero
section.

\subsection*{Proposition 3.3.}
If \(X\) is represented by an \(S\)-scheme \(\widetilde X\), then
\(T_{X/S}(\mathcal M)\) and \(L^u_{X/S}(\mathcal M)\) are representable.
In particular, \(T_{X/S}\) and \(L^u_{X/S}\) are represented by the
vector fibrations
\[
  V(\Omega^1_{\widetilde X/S})
  \qquad\text{and}\qquad
  V(u^*\Omega^1_{\widetilde X/S})
\]
over \(\widetilde X\) and \(S\), respectively.

It plainly suffices to prove the assertion for \(T_{X/S}(\mathcal M)\),
since the corresponding assertion for \(L^u_{X/S}(\mathcal M)\) follows
by pullback. By Corollary~2.2.1, it even suffices to treat \(T_{X/S}\),
in which case the result is exactly Nota~2.2.2.

% Source page: Exp2 component-local p. 9; combined-reader p. 67.
\subsection*{Remark 3.3.1.}
The proposition gives a particularly simple description of the vector
fibration representing \(L^u_{X/S}\). The image of the section \(u\) of
\(X\) over \(S\) is locally closed,
\footnote{Editor's note: cf. EGA~I, 5.3.11.}
and hence is defined by a quasi-coherent ideal \(\mathfrak m\) in the
scheme induced on an open subscheme of \(X\). The quotient
\(\mathfrak m/\mathfrak m^2\) may be regarded as a quasi-coherent module
on \(S\); it is this module that defines the vector fibration in question.

For example, let \(X\) be an algebraic scheme over a field \(k\), and let
\(u\) be a \(k\)-rational point of \(X\). Let \(\mathfrak m\) be the
maximal ideal of the local ring \(\mathcal O_{X,u}\), and let
\(\mathfrak t\) be the \(k\)-vector space dual to
\(\mathfrak m/\mathfrak m^2\); this is the Zariski tangent space of
\(\mathcal O_{X,u}\) at \(u\). In the notation of I~4.6.5.1,
\[
  L^u_{X/k}=V(\mathfrak m/\mathfrak m^2)=W(\mathfrak t).
\]

Returning to the general situation, observe first that
\(L^u_{X/S}(\mathcal M)\) is covariant from finite free
\(\mathcal O_S\)-modules to functors over \(S\). In particular,
\(\mathcal O(S)\) is a set of operators of the \(S\)-functor
\(L^u_{X/S}(\mathcal M)\), compatible with functoriality in
\(\mathcal M\).
\footnote{Editor's note: cf. 3.1.2.}

\subsection*{Proposition 3.4.}
Formation of \(T_{X/S}(\mathcal M)\) and \(L^u_{X/S}(\mathcal M)\)
commutes with base extension. For every \(S\)-scheme \(S'\), there are
isomorphisms, functorial in \(\mathcal M\),
\[
  T_{X_{S'}/S'}(\mathcal M\otimes_{\mathcal O_S}\mathcal O_{S'})
    \xrightarrow{\sim} T_{X/S}(\mathcal M)_{S'},
\]
\[
  L^{u'}_{X_{S'}/S'}(\mathcal M\otimes_{\mathcal O_S}\mathcal O_{S'})
    \xrightarrow{\sim} L^u_{X/S}(\mathcal M)_{S'},
  \qquad u'=u_{S'}.
\]
This follows immediately from the compatibility of
\(\underline{\operatorname{Hom}}\) with base extension.

\subsection*{Corollary 3.4.1.}
The \(X\)-functor \(T_{X/S}(\mathcal M)\), respectively the \(S\)-functor
\(L^u_{X/S}(\mathcal M)\), is naturally an object with operators
\(\mathcal O_X\), respectively \(\mathcal O_S\), and these structures
are functorial in \(\mathcal M\).

For \(L^u_{X/S}(\mathcal M)\), this is immediate. For each \(S'\) over
\(S\), the ring \(\mathcal O(S')\) acts on
\(\mathcal M\otimes\mathcal O_{S'}\), hence on
\[
  L^{u'}_{X_{S'}/S'}(\mathcal M\otimes\mathcal O_{S'})
  =L^u_{X/S}(\mathcal M)_{S'};
\]
this action is functorial in \(S'\).

The case of \(T_{X/S}(\mathcal M)\) requires a little more detail. For an
\(X\)-object \(f:X'\to X\), put
\(T_{X/S}(\mathcal M)_{X'}=T_{X/S}(\mathcal M)\times_X X'\). We must
equip
\[
  T_{X/S}(\mathcal M)_{X'}(X')
  =\operatorname{Hom}_X\bigl(X',T_{X/S}(\mathcal M)\bigr)
\]
with an action of \(\mathcal O(X')\), functorially in \(X'\). Let
\(X_{X'}=X\times_S X'\), and let \(f'\) be the section of \(X_{X'}\)
over \(X'\) defined by \(f\). The following diagram provides the required
identification.

% Source page: Exp2 component-local p. 10; combined-reader p. 68.
\SGAThreeDiagram[0.39\textwidth]
  {figures/exp2/exp2_p010_base_change_operator_diagram.png}
  {The base-change diagram used to define the \(\mathcal O_X\)-operator structure.}
  {fig:exp2-base-change-operator}
  {Polo--Gille Exp2-14oct24.pdf, component-local p. 10; combined-reader p. 68}

Together with 3.2.1, the diagram identifies
\(T_{X/S}(\mathcal M)_{X'}(X')\) with
\[
  L^{f'}_{X_{X'}/X'}(\mathcal M)(X')
  =\left\{
    \begin{array}{c|c}
      X'\text{-morphisms }\psi:I_{X'}(\mathcal M)\to X_{X'}
      &\psi\circ\varepsilon_{\mathcal M}=f'
    \end{array}
  \right\}.
\]
Every \(a\in\mathcal O(X')\) acts through its action on
\(I_{X'}(\mathcal M)\): in the notation of 2.1.1,
\[
  a\psi=\psi\circ a^*,
  \qquad (a\psi)(x)=\psi(x\cdot a)
\]
for every \(X''\to X'\) and every \(x\in I_{X'}(\mathcal M)(X'')\).
This construction is plainly functorial in \(X'\).
\footnote{Editor's note: the original has been expanded in what follows.}
The isomorphisms of Proposition~3.4 are therefore, by construction,
isomorphisms for the corresponding \(\mathcal O_{X_{S'}}\)- and
\(\mathcal O_{S'}\)-operator structures.

\subsection*{Remark 3.4.2.}
\footnote{Editor's note: Remarks 3.4.2 and 3.4.3 have been added.}
The action of \(\mathcal O_X\) on \(T_{X/S}(\mathcal M)\) can be seen more
directly as follows. For every \(f:X'\to X\),
\[
  \operatorname{Hom}_X\bigl(X',T_{X/S}(\mathcal M)\bigr)
  =\left\{\varphi\in
    \operatorname{Hom}_S(I_{X'}(\mathcal M),X)
    \ \middle|\ \varphi\circ\varepsilon_{\mathcal M}=f\right\}.
\]
By 2.1.3, the \(S\)-functor \(I_{X'}(\mathcal M)\) carries an action of
the monoid \(\mathcal O(X')\) preserving the zero section. Thus, if
\(a^*\) denotes the \(X'\)-endomorphism defined by
\(a\in\mathcal O(X')\), then
\[
  a\varphi=\varphi\circ a^*,
  \qquad
  (a\varphi)(m,x')=\varphi(m\cdot a(x'),x')
\]
for every \(S'\to S\) and every
\((m,x')\in\operatorname{Hom}_S
  (S',I_S(\mathcal M)\times_S X')\). Notice that
\(a^*\circ\varepsilon_{\mathcal M}=\varepsilon_{\mathcal M}\), so that
\((a\varphi)\circ\varepsilon_{\mathcal M}=f\).

Likewise, for \(t:S'\to S\), the set
\(L^u_{X/S}(\mathcal M)(S')\) consists of the \(S\)-morphisms
\(\varphi:I_{S'}(\mathcal M)\to X\) satisfying
\(\varphi\circ\varepsilon_{\mathcal M}=u\circ t\), and
\(a\varphi=\varphi\circ a^*\) for \(a\in\mathcal O(S')\).

\subsection*{Remark 3.4.3.}
\footnotemark[21]
The observations of 3.1.2 apply equally to the action of \(\mathcal O_S\)
on \(L^u_{X/S}(\mathcal M)\) and to that of \(\mathcal O_X\) on
\(T_{X/S}(\mathcal M)\).

% Source page: Exp2 component-local p. 11; combined-reader p. 69.
\subsection*{Definition 3.5.}
Let \(S\) be a scheme and \(X\) an \(S\)-functor. We say that \(X\)
\emph{satisfies condition (E) relative to \(S\)} if, for every \(S'\)
over \(S\) and every pair of finite free \(\mathcal O_{S'}\)-modules
\(\mathcal M,\mathcal N\), the diagram of sets obtained by applying \(X\)
to diagram \((*)\) of 2.1 is cartesian.
\footnote{Editor's note: for examples of \(S\)-functors and \(S\)-groups
that do not satisfy (E), see 6.2 and 6.3.}

\SGAThreeDiagram[0.46\textwidth]
  {figures/exp2/exp2_p011_condition_E_cartesian_diamond.png}
  {The cartesian diagram in condition (E).}
  {fig:exp2-condition-e}
  {Polo--Gille Exp2-14oct24.pdf, component-local p. 11; combined-reader p. 69}

\subsection*{3.5.1.}
\footnote{Editor's note: what follows has been added.}
Equivalently, the functor
\(\mathcal M\mapsto T_{X/S}(\mathcal M)\) carries direct sums of finite
free \(\mathcal O_S\)-modules to products of \(X\)-functors. In that case,
the same is true of
\[
  \mathcal M\longmapsto L^u_{X/S}(\mathcal M)
  =S\times_XT_{X/S}(\mathcal M)
\]
for every \(u\in\Gamma(X/S)\).

Proposition~2.2 shows that every representable functor satisfies
condition (E). We shall sometimes abbreviate ``\(X\) satisfies condition
(E) relative to \(S\)'' to ``\(X/S\) satisfies condition (E).''

If \(X/S\) satisfies (E), then
\(\mathcal M\mapsto T_{X/S}(\mathcal M)\) commutes with products and hence
carries groups to groups. In particular, \(T_{X/S}(\mathcal M)\) is a
commutative \(X\)-group; for the same reason,
\(L^u_{X/S}(\mathcal M)\) is a commutative \(S\)-group.

\subsection*{Proposition 3.6.}
If \(X/S\) satisfies (E), then the abelian-group structure on
\(T_{X/S}(\mathcal M)\), respectively on \(L^u_{X/S}(\mathcal M)\),
together with the action of \(\mathcal O_X\), respectively of
\(\mathcal O_S\), makes it an \(\mathcal O_X\)-module, respectively an
\(\mathcal O_S\)-module.

The operator action is functorial in \(\mathcal M\), and therefore
respects the abelian-group structure obtained by functoriality from that
of \(\mathcal M\).
\footnote{Editor's note: what follows has been expanded.}
Indeed, retaining the notation of 3.1.2, the abelian \(X\)-group law on
\(T_{X/S}(\mathcal M)\) is the composite
\[
  T_{X/S}(\mathcal M)\times_XT_{X/S}(\mathcal M)
  \simeq T_{X/S}(\mathcal M\oplus\mathcal M)
  \xrightarrow{\ m\ }T_{X/S}(\mathcal M),
\]
whereas
\[
  T_{X/S}(\mathcal M)\xrightarrow{\ \delta\ }
  T_{X/S}(\mathcal M\oplus\mathcal M)
  \simeq T_{X/S}(\mathcal M)\times_XT_{X/S}(\mathcal M)
\]
is the diagonal morphism.

% Source page: Exp2 component-local p. 12; combined-reader p. 70.
By Remark~3.4.3 and the identities \((\dagger)\) of 3.1.2,
\[
  \lambda(x+y)=\lambda x+\lambda y,
  \qquad
  (\lambda+\mu)x=\lambda x+\mu x
\]
for every \(f:X'\to X\), every
\(x,y\in\operatorname{Hom}_X(X',T_{X/S}(\mathcal M))\), and every
\(\lambda,\mu\in\mathcal O(X')\).

\subsection*{Remark 3.6.1.}
If \(X\) is representable, then it satisfies (E), and \(T_{X/S}\) and
\(L^u_{X/S}\) are represented by vector fibrations. The preceding laws
are then exactly those induced by the vector-fibration structures
(cf. I~4.6).
\footnote{Editor's note: explicitly, for every \(S\)-morphism
\(f:X'\to X\), the action of \(\mathcal O(X')\) on
\(\operatorname{Hom}_X(X',T_{X/S}(\mathcal M))\) corresponds, under the
identification
\[
  \operatorname{Hom}_X(X',T_{X/S}(\mathcal M))
  \simeq
  \operatorname{Hom}_{\mathcal O_{X'}}
    (f^*\Omega^1_{X/S},\mathcal M\otimes_{\mathcal O_S}\mathcal O_{X'}),
\]
to the natural action on the right. This follows from the proof of SGA~1,
III~5.1; see also the addition 0.1.8 in Expos\'e~III.}

\subsection*{Proposition 3.4 bis.}
If \(X/S\) satisfies (E), then \(X_{S'}/S'\) satisfies (E), and the
isomorphisms of 3.4 respect the \(\mathcal O_{X_{S'}}\)-module and
\(\mathcal O_{S'}\)-module structures, respectively.

No further comment is needed.

\subsection*{Proposition 3.7.}
The functors \(T_{X/S}(\mathcal M)\) and \(L^u_{X/S}(\mathcal M)\) are
functorial in \(X\). Thus, if \(f:X\to X'\) is an \(S\)-morphism, the
following diagrams commute.

\SGAThreeDiagram[0.74\textwidth]
  {figures/exp2/exp2_p012_functoriality_T_L_diagrams.png}
  {Functoriality of tangent bundles and tangent spaces in \(X\).}
  {fig:exp2-tangent-functoriality}
  {Polo--Gille Exp2-14oct24.pdf, component-local p. 12; combined-reader p. 70}

Moreover, if \(f\) is a monomorphism, then so are \(T(f)\) and \(L(f)\).
\footnote{Editor's note: the preceding sentence has been added.}
The existence of \(T(f)\) and \(L(f)\), and the last assertion, follow
immediately from the definitions. Commutativity follows from functoriality
in \(\mathcal M\) and from \(T_{X/S}(0)=X\).

\subsection*{Remark 3.7.1.}
\footnote{Editor's note: the hypothesis that \(X\) and \(X'\) be
representable has been added, and what follows has been expanded.}
Suppose that \(X\) and \(X'\) are representable, and let \(r\) be the rank
of \(\mathcal M\). By 2.2.2, \(T_{X/S}(\mathcal M)\) is isomorphic to the
product over \(X\) of \(r\) copies of \(V(\Omega^1_{X/S})\), and similarly
for \(T_{X'/S}(\mathcal M)\). Hence the square in 3.7 is cartesian when
\(f\) is an open immersion, or more generally when
\(f^*\Omega^1_{X'/S}=\Omega^1_{X/S}\), for example when \(f\) is
\'etale. Under these conditions there is an isomorphism of \(S\)-functors
\[
  L^u_{X/S}(\mathcal M)
  \xrightarrow{\sim}L^{f\circ u}_{X'/S}(\mathcal M).
\]

% Source page: Exp2 component-local p. 13; combined-reader p. 71.
In general, the square in 3.7 defines the following morphism of
\(X\)-functors.

\SGAThreeDiagram[0.42\textwidth]
  {figures/exp2/exp2_p013_general_X_functor_triangle.png}
  {The induced morphism over \(X\) in the general case.}
  {fig:exp2-general-x-functor}
  {Polo--Gille Exp2-14oct24.pdf, component-local p. 13; combined-reader p. 71}

\subsection*{Proposition 3.7 bis.}
Let \(f:X\to X'\) be an \(S\)-morphism. If \(X\) and \(X'\) satisfy
(E) relative to \(S\), then
\[
  T_{X/S}(\mathcal M)\xrightarrow{\ T(f)\ }
  T_{X'/S}(\mathcal M)_X,
  \qquad\text{respectively}\qquad
  L^u_{X/S}(\mathcal M)\xrightarrow{\ L(f)\ }
  L^{f\circ u}_{X'/S}(\mathcal M),
\]
is a morphism of \(\mathcal O_X\)-modules, respectively of
\(\mathcal O_S\)-modules.

This follows from Proposition~3.7 by functoriality in \(\mathcal M\).

\subsection*{Proposition 3.8.}
For two functors \(X,Y\) over \(S\), there are isomorphisms, functorial in
\(\mathcal M\),
\begin{align}
  T_{X/S}(\mathcal M)\times_ST_{Y/S}(\mathcal M)
    &\xrightarrow{\sim}T_{(X\times_SY)/S}(\mathcal M), \tag{3.8.1}\\
  L^u_{X/S}(\mathcal M)\times_SL^v_{Y/S}(\mathcal M)
    &\xrightarrow{\sim}L^{(u,v)}_{(X\times_SY)/S}(\mathcal M). \tag{3.8.2}
\end{align}
The first isomorphism follows from 1.2 \((*)\);
\footnote{Editor's note: what follows has been expanded.}
the second follows from it
by base change along \((u,v):S\to X\times_SY\).

\subsection*{Remark 3.8.0.}
The isomorphism (3.8.1) may also be read as an isomorphism of
\(X\times_SY\)-functors
\[
  \bigl(T_{X/S}(\mathcal M)\times_X(X\times_SY)\bigr)
  \mathop{\times}_{X\times_SY}
  \bigl(T_{Y/S}(\mathcal M)\times_Y(X\times_SY)\bigr)
  \xrightarrow{\sim}T_{(X\times_SY)/S}(\mathcal M).
\]

\subsection*{Corollary 3.8.1.}
If \(X/S\) carries an algebraic structure defined by finite cartesian
products, then \(T_{X/S}(\mathcal M)\) carries a structure of the same
kind, and the projection \(T_{X/S}(\mathcal M)\to X\) is a morphism for
that structure.

\subsection*{Proposition 3.8 bis.}
If \(X/S\) and \(Y/S\) satisfy (E), then \((X\times_SY)/S\) satisfies
(E), and (3.8.1), respectively (3.8.2), is an isomorphism of
\(\mathcal O_{X\times_SY}\)-modules, respectively of
\(\mathcal O_S\)-modules.

\emph{Proof.}
\footnote{Editor's note: this proof has been added.}
By 3.5.1 and (3.8.1), the product \((X\times_SY)/S\) satisfies (E). Let
us prove the assertion about (3.8.1). Given an \(S\)-morphism
\((x,y):Z\to X\times_SY\), Remark~3.4.2 reduces the claim to the
\(\mathcal O(Z)\)-linearity of the map
\begin{multline*}
  \{\varphi\in\operatorname{Hom}_S(I_Z(\mathcal M),X)
       \mid\varphi\circ\varepsilon_{\mathcal M}=x\}
  \times
  \{\psi\in\operatorname{Hom}_S(I_Z(\mathcal M),Y)
       \mid\psi\circ\varepsilon_{\mathcal M}=y\}\\
  \longrightarrow
  \{\theta\in\operatorname{Hom}_S(I_Z(\mathcal M),X\times_SY)
       \mid\theta\circ\varepsilon_{\mathcal M}=(x,y)\},
\end{multline*}
which sends \((\varphi,\psi)\) to \(\varphi\times\psi\). This is clear,
since for \(a\in\mathcal O(Z)\),
\[
  (\varphi\circ a^*)\times(\psi\circ a^*)
  =(\varphi\times\psi)\circ a^*
  =a\cdot(\varphi\times\psi).
\]
The assertion about (3.8.2) follows similarly from 3.2.1.

% Source page: Exp2 component-local p. 14; combined-reader p. 72.
\subsection*{3.9.0.}
\footnote{Editor's note: this paragraph has been added to explain the
introduction of the notion of an \(S\)-\(H\)-functor.}
If \(X\) is an \(S\)-group with unit section \(e:S\to X\), write
\[
  \underline{\operatorname{Lie}}(X/S,\mathcal M)
  =L^e_{X/S}(\mathcal M).
\]
Thus \(\underline{\operatorname{Lie}}(X/S,\mathcal M)\) is defined by the
following cartesian square.

\SGAThreeDiagram[0.38\textwidth]
  {figures/exp2/exp2_p014_lie_kernel_cartesian_square.png}
  {The cartesian square defining \(\underline{\operatorname{Lie}}(X/S,\mathcal M)\).}
  {fig:exp2-lie-kernel}
  {Polo--Gille Exp2-14oct24.pdf, component-local p. 14; combined-reader p. 72}

By Corollary~3.8.1, the projection
\(p:T_{X/S}(\mathcal M)\to X\) is a morphism of \(S\)-groups. Hence
\(\underline{\operatorname{Lie}}(X/S,\mathcal M)\) is an \(S\)-group,
identified by \(i\) with the kernel of \(p\).

If \(X/S\) also satisfies (E), Proposition~3.9 will show that this group
structure, induced by that of \(X\), agrees with the abelian-group
structure induced by functoriality in \(\mathcal M\) (cf. 3.5.1). In fact,
the result holds under the weaker assumption that \(X\) be an
\(S\)-functor in monoids, or more generally an \(S\)-functor in
\(H\)-sets, as defined next.

\subsection*{Definitions 3.9.0.1.}
\begin{enumerate}
\item[(a)] An \emph{\(H\)-set} is a set \(X\) equipped with a composition
law having a two-sided identity, denoted \(e_X\) or simply \(e\).
\footnote{Editor's note: this terminology is inspired by the notion of an
\(H\)-space in topology.}
If \(f:X\to Y\) is a morphism of \(H\)-sets, its kernel is
\(\ker f=f^{-1}(e_Y)\), a sub-\(H\)-set of \(X\).

\item[(b)] An \emph{\(H\)-object} in a category \(\mathscr C\) is an
object \(X\) equipped with a morphism \(X\times X\to X\) and a section
of \(X\) over the final object that is a two-sided identity. Every
\(\mathscr C\)-monoid, and in particular every \(\mathscr C\)-group, is
a \(\mathscr C\)-\(H\)-object. An \(H\)-object in the category of functors
over \(S\) will be called an \emph{\(S\)-\(H\)-functor}.

% Source page: Exp2 component-local p. 15; combined-reader p. 73.
\item[(c)] If \(X\) is an \(S\)-\(H\)-functor (for example, an
\(S\)-group), and \(e:S\to X\) is its identity section, write
\[
  \underline{\operatorname{Lie}}(X/S,\mathcal M)=L^e_{X/S}(\mathcal M),
  \qquad
  \underline{\operatorname{Lie}}(X/S)
  =\underline{\operatorname{Lie}}(X/S,\mathcal O_S).
\]
\end{enumerate}

We now spell out the following special case of 3.8.1.
\footnote{Editor's note: Corollary 3.9.0.2 has been added; it will be used
in 4.1.}

\subsection*{Corollary 3.9.0.2.}
If \(X\) is an \(S\)-\(H\)-functor, respectively an \(S\)-group, then
\(T_{X/S}(\mathcal M)\) and
\(\underline{\operatorname{Lie}}(X/S,\mathcal M)\) are likewise
\(S\)-\(H\)-functors, respectively \(S\)-groups, and there are morphisms
of \(S\)-\(H\)-functors, respectively of \(S\)-groups, as follows.

\SGAThreeDiagram[0.50\textwidth]
  {figures/exp2/exp2_p015_lie_tangent_split_sequence.png}
  {The inclusion, projection, and zero section for the tangent functor.}
  {fig:exp2-lie-tangent-split}
  {Polo--Gille Exp2-14oct24.pdf, component-local p. 15; combined-reader p. 73}

Here \(i\) identifies \(\underline{\operatorname{Lie}}(X/S,\mathcal M)\)
with \(\ker p\), and \(s\) is a section of \(p\).

\subsection*{Proposition 3.9.}
Let \(X\) be an \(S\)-\(H\)-functor satisfying (E) relative to \(S\).
The \(S\)-\(H\)-functor structure on
\(\underline{\operatorname{Lie}}(X/S,\mathcal M)\) induced by that of
\(X\) agrees with the \(S\)-group structure defined in 3.5.1.

The preceding discussion shows that
\(\underline{\operatorname{Lie}}(X/S,\mathcal M)\) is an \(H\)-object in
the category of \(\mathcal O_S\)-modules. The proposition follows from
the next lemma.
\footnote{Editor's note: this is inspired by the standard proof that the
fundamental group of an \(H\)-space is abelian.}

\subsection*{Lemma 3.10.}
Let \(\mathscr C\) be a category. Let \(G\) be an \(H\)-object in the
category of \(\mathscr C\)-\(H\)-objects. Thus \(G\) is a
\(\mathscr C\)-\(H\)-object, whose composition law we denote
\(f:G\times G\to G\), together with a morphism of
\(\mathscr C\)-\(H\)-objects \(h:G\times G\to G\).
\footnote{Editor's note: \(G\times G\) is equipped with the law
\((G\times G)\times(G\times G)\to G\times G\),
\(((x,z),(y,t))\mapsto(f(x,y),f(z,t))\).}
Then \(f=h\), and \(f\) is commutative.

Evaluating the functors on a variable argument reduces the proof, as
usual, to the case in which \(\mathscr C\) is the category of sets. We
then have a set \(G\) and two maps \(f,h:G\times G\to G\) satisfying
\[
  h(f(x,y),f(z,t))=f(h(x,z),h(y,t)).
\]
Let \(e,u\in G\) be the respective identities for \(f\) and \(h\), so
that
\[
  f(e,x)=f(x,e)=x,
  \qquad h(u,x)=h(x,u)=x.
\]
First,
\[
  h(f(u,y),f(x,u))=f(x,y)=h(f(x,u),f(u,y)).
\]
Putting \(y=e\), respectively \(x=e\), gives
\begin{align*}
  x&=f(x,e)=h(f(u,e),f(x,u))=h(u,f(x,u))=f(x,u),\\
  y&=f(e,y)=h(f(e,u),f(u,y))=h(u,f(u,y))=f(u,y).
\end{align*}
Substituting these identities into the original relation yields
\[
  h(y,x)=f(x,y)=h(x,y).
\]
This proves the lemma and Proposition~3.9. The following corollaries now
follow from 3.9.

\subsection*{Corollary 3.9.1.}
If \(X\) is an \(S\)-\(H\)-functor satisfying (E) relative to \(S\),
every element of \(X(I_S(\mathcal M))\) that maps to the identity element
of \(X(S)\) is invertible.

% Source page: Exp2 component-local p. 16; combined-reader p. 74.
\subsection*{Corollary 3.9.2.}
If \(X\) is an \(S\)-monoid satisfying (E) relative to \(S\), an element
of \(X(I_S(\mathcal M))\) is invertible if and only if its image in
\(X(S)\) is invertible.

\subsection*{Corollary 3.9.3.}
If \(X\) is an \(S\)-group satisfying (E) relative to \(S\), the two
\(S\)-group laws on
\(\underline{\operatorname{Lie}}(X/S,\mathcal M)\) coincide.

\subsection*{Corollary 3.9.4.}
\footnote{Editor's note: this corollary has been added; it amplifies
Remark 4.1.1.2 below.}
Let \(G\) be an \(S\)-group satisfying (E) relative to \(S\). For
\(n\in\mathbb Z\), let \(n_G:G\to G\) be the morphism of \(S\)-functors
defined by \(g\mapsto g^n\). Then the derived morphism
\[
  L(n_G):\underline{\operatorname{Lie}}(G/S)\longrightarrow
  \underline{\operatorname{Lie}}(G/S)
\]
is multiplication by \(n\); that is, it sends
\(x\in\underline{\operatorname{Lie}}(G/S)(S')\) to \(nx\).

The map \(n_G\) is not generally a group homomorphism, but it preserves
the identity section \(e:S\to G\), so its derivative indeed carries
\(\underline{\operatorname{Lie}}(G/S)=L^e_{G/S}\) to itself. If \(i\)
denotes the inclusion
\(\underline{\operatorname{Lie}}(G/S)\hookrightarrow T_{G/S}\), then
\(L(n_G)\) is characterized by
\[
  i(L(n_G)(x))=i(x)^n.
\]
By 3.9, the group laws on \(\underline{\operatorname{Lie}}(G/S)\) coming
from (E) and from the law on \(G\) agree. Hence
\(i(x)^n=i(nx)\), and therefore \(L(n_G)(x)=nx\).

Before drawing further consequences from Proposition~3.9, we prove one
more functoriality result.

\subsection*{Proposition 3.11.}
In the situation of Section~1, there is an isomorphism, functorial in
\(\mathcal M\),
\[
  T_{\underline{\operatorname{Hom}}_{Z/S}(X,Y)/S}(\mathcal M)
  \xrightarrow{\sim}
  \underline{\operatorname{Hom}}_{Z/S}
    \bigl(X,T_{Y/Z}(\mathcal M)\bigr).
\]
Indeed, by definition,
\[
  T_{\underline{\operatorname{Hom}}_{Z/S}(X,Y)/S}(\mathcal M)
  =\underline{\operatorname{Hom}}_S
    \bigl(I_S(\mathcal M),\underline{\operatorname{Hom}}_{Z/S}(X,Y)\bigr).
\]
Applying the isomorphism (4) of 1.1 with \(T=I_S(\mathcal M)\) gives
\begin{multline*}
  \underline{\operatorname{Hom}}_S
    \bigl(I_S(\mathcal M),\underline{\operatorname{Hom}}_{Z/S}(X,Y)\bigr)\\
  \simeq
  \underline{\operatorname{Hom}}_{Z/S}
    \bigl(X,\underline{\operatorname{Hom}}_Z
      (Z\times_SI_S(\mathcal M),Y)\bigr).
\end{multline*}
Since \(Z\times_SI_S(\mathcal M)\simeq I_Z(\mathcal M)\), this is
\[
  \underline{\operatorname{Hom}}_{Z/S}
    \bigl(X,\underline{\operatorname{Hom}}_Z(I_Z(\mathcal M),Y)\bigr)
  =\underline{\operatorname{Hom}}_{Z/S}
    \bigl(X,T_{Y/Z}(\mathcal M)\bigr).
\]

\subsection*{Corollary 3.11.1.}
If \(Y/Z\) satisfies (E), then
\(\underline{\operatorname{Hom}}_{Z/S}(X,Y)/S\) satisfies (E), and the
isomorphism of 3.11 respects the module structures over
\(\underline{\operatorname{Hom}}_{Z/S}(X,Y)\).
\footnote{Editor's note: writing
\(H=\underline{\operatorname{Hom}}_{Z/S}(X,Y)\), this implies in
particular that \(T_{H/S}(\mathcal M)\) is naturally a module over
\(\prod_{X\times_SH/H}\mathcal O_H\). This result appeared somewhat
surprising to the editors, so the proof has been given in detail.}

\emph{Proof.}
Let \(\mathcal M,\mathcal N\) be finite free \(\mathcal O_S\)-modules.
If \(Y/Z\) satisfies (E), then
\[
  T_{Y/Z}(\mathcal M\oplus\mathcal N)
  \simeq T_{Y/Z}(\mathcal M)\times_YT_{Y/Z}(\mathcal N).
\]
The right-hand side is a subfunctor of
\(T_{Y/Z}(\mathcal M)\times_ST_{Y/Z}(\mathcal N)\), and 1.2 \((*)\)
therefore gives
\begin{multline*}
  \underline{\operatorname{Hom}}_{Z/S}
    \bigl(X,T_{Y/Z}(\mathcal M\oplus\mathcal N)\bigr)\\
  \simeq
  \underline{\operatorname{Hom}}_{Z/S}\bigl(X,T_{Y/Z}(\mathcal M)\bigr)
  \mathop{\times}_{\underline{\operatorname{Hom}}_{Z/S}(X,Y)}
  \underline{\operatorname{Hom}}_{Z/S}\bigl(X,T_{Y/Z}(\mathcal N)\bigr).
\end{multline*}
Together with 3.11, this says that
\(\underline{\operatorname{Hom}}_{Z/S}(X,Y)\) satisfies (E), proving the
first assertion.

For the second, put
\(H=\underline{\operatorname{Hom}}_{Z/S}(X,Y)\), and let
\(\Delta:H'\to H\) be an \(S\)-morphism. Equivalently, let
\(\delta:H'\times_SX\to Y\) be the corresponding \(Z\)-morphism; it
makes \(H'\times_SX\) a \(Y\)-object. On the one hand, we have the
following commutative diagram.

% Source page: Exp2 component-local p. 17; combined-reader p. 75.
\SGAThreeDiagram[0.82\textwidth]
  {figures/exp2/exp2_p017_hom_module_comparison_diagram.png}
  {The first comparison diagram in the proof of Corollary 3.11.1.}
  {fig:exp2-hom-module-comparison}
  {Polo--Gille Exp2-14oct24.pdf, component-local p. 17; combined-reader p. 75}

By 1.3, the action of \(\alpha\in\mathcal O(H'\times_SX)\) on
\(\Psi\in\operatorname{Hom}_Y(H'\times_SX,T_{Y/Z}(\mathcal M))\) is
given, for every \(U\to S\) and every
\((h,x)\in\operatorname{Hom}_S(U,H'\times_SX)\), by
\[
  (\alpha\Psi)(h,x)=\alpha(h,x)\Psi(h,x),
\]
where \(U\) is viewed over \(Y\) through \(\delta\circ(h,x)\), and
\(\alpha(h,x)\in\mathcal O(U)\) acts through the
\(\mathcal O_Y\)-module structure on \(T_{Y/Z}(\mathcal M)\). By 3.4.2,
under the identification
\[
  \operatorname{Hom}_Y(H'\times_SX,T_{Y/Z}(\mathcal M))
  =\{\psi\in\operatorname{Hom}_Z(I_{H'\times_SX}(\mathcal M),Y)
       \mid\psi\circ\varepsilon_{\mathcal M}=\delta\},
\]
this action is
\begin{equation}
  (\alpha\psi)(m,h,x)
  =\psi(m\cdot\alpha(h,x),h,x)
  \tag{1}
\end{equation}
for every
\((m,h,x)\in\operatorname{Hom}_S
  (U,I_S(\mathcal M)\times_SH'\times_SX)\).

% Source page: Exp2 component-local p. 18; combined-reader p. 76.
On the other hand, consider
\(T_{H/S}(\mathcal M)=\underline{\operatorname{Hom}}_S
  (I_S(\mathcal M),H)\). We have a second commutative diagram.

\SGAThreeDiagram[0.84\textwidth]
  {figures/exp2/exp2_p018_tangent_hom_comparison_diagram.png}
  {The second comparison diagram in the proof of Corollary 3.11.1.}
  {fig:exp2-tangent-hom-comparison}
  {Polo--Gille Exp2-14oct24.pdf, component-local p. 18; combined-reader p. 76}

The bijection \((*)\) in the diagram is described as follows (cf. 1.1(2)
and I~1.7.1). For every \(U\to S\) and
\((m,h,x)\in\operatorname{Hom}_S
  (U,I_S(\mathcal M)\times_SH'\times_SX)\), where \(U\) lies over \(Z\)
through \(U\xrightarrow{x}X\to Z\), one has
\(\Phi(m,h)\in\operatorname{Hom}_Z(X\times_SU,Y)\) and
\begin{equation}
  \varphi(m,h,x)=\Phi(m,h)\circ(x\times\operatorname{id}_U)
  \in\operatorname{Hom}_Z(U,Y).
  \tag{$\dagger$}
\end{equation}
By 3.4.2, with \(X\) replaced by \(H\) and \(X'\) by \(H'\), the
action of \(a\in\mathcal O(H')\) on
\(\Phi\in\operatorname{Hom}_S(I_{H'}(\mathcal M),H)\) is
\[
  (a\Phi)(m,h)=\Phi(m\cdot a(h),h).
\]
Consequently, if \(\varphi\), respectively \(a\varphi\), corresponds to
\(\Phi\), respectively \(a\Phi\), then \((\dagger)\) gives
\begin{equation}
  (a\varphi)(m,h,x)
  =\Phi(m\cdot a(h),h)\circ(x\times\operatorname{id}_U)
  =\varphi(m\cdot a(h),h,x).
  \tag{2}
\end{equation}
Together with (1), this proves that the isomorphism
\[
  T_{H/S}(\mathcal M)
  \xrightarrow{\sim}
  \underline{\operatorname{Hom}}_{Z/S}
    (X,T_{Y/Z}(\mathcal M))
\]
of 3.11.1 is an isomorphism of \(\mathcal O(H)\)-modules. Moreover, for
every \(H'\to H\), the \(\mathcal O(H')\)-module structure on
\(\operatorname{Hom}_H(H',T_{H/S}(\mathcal M))\) extends, functorially in
\(H'\), to an \(\mathcal O(H'\times_SX)\)-module structure.

Taking \(Z=S\) yields the next corollary.

\subsection*{Corollary 3.11.2.}
There is an isomorphism, functorial in \(\mathcal M\),
\[
  T_{\underline{\operatorname{Hom}}_S(X,Y)/S}(\mathcal M)
  \xrightarrow{\sim}
  \underline{\operatorname{Hom}}_S(X,T_{Y/S}(\mathcal M)).
\]
If \(Y/S\) satisfies (E), then
\(\underline{\operatorname{Hom}}_S(X,Y)/S\) satisfies (E), and this
isomorphism respects the module structures over
\(\underline{\operatorname{Hom}}_S(X,Y)\).

Let \(u:X\to Y\) be an \(S\)-morphism, identified with the constant
morphism \(u:S\to\underline{\operatorname{Hom}}_S(X,Y)\) such that
\(u(f)=u\) for every \(f:S'\to S\). The fiber product of \(u\) with
\[
  \underline{\operatorname{Hom}}_S(X,T_{Y/S}(\mathcal M))
  \longrightarrow\underline{\operatorname{Hom}}_S(X,Y)
\]
identifies with
\(\underline{\operatorname{Hom}}_{Y/S}(X,T_{Y/S}(\mathcal M))\), where
\(X\) lies over \(Y\) through \(u\). Hence the definition of the tangent
space and Corollary~3.11.2 give the following result.
\footnote{Editor's note: the preceding sentences have been added.}

% Source page: Exp2 component-local p. 19; combined-reader p. 77.
\subsection*{Corollary 3.11.3.}
For an \(S\)-morphism \(u:X\to Y\), there is an isomorphism, functorial
in \(\mathcal M\),
\[
  L^u_{\underline{\operatorname{Hom}}_S(X,Y)/S}(\mathcal M)
  \xrightarrow{\sim}
  \underline{\operatorname{Hom}}_{Y/S}(X,T_{Y/S}(\mathcal M)),
\]
where on the right \(X\) lies over \(Y\) through \(u\). If \(Y/S\)
satisfies (E), this is an isomorphism of \(\mathcal O_S\)-modules.
\footnote{Editor's note: the preceding sentence has been added.}

In particular, for \(Y=X\), the functor
\(\underline{\operatorname{End}}_S(X)\) is an \(S\)-functor in monoids,
hence an \(S\)-\(H\)-functor. Recall that
\(\underline{\operatorname{Lie}}(\underline{\operatorname{End}}_S(X)/S,
\mathcal M)\) denotes
\(L^e_{\underline{\operatorname{End}}_S(X)/S}(\mathcal M)\), with \(e\)
the identity section.

\subsection*{Corollary 3.11.4.}
There is an isomorphism, functorial in \(\mathcal M\),
\[
  \underline{\operatorname{Lie}}
    (\underline{\operatorname{End}}_S(X)/S,\mathcal M)
  \xrightarrow{\sim}
  \prod_{X/S}T_{X/S}(\mathcal M).
\]
If \(X/S\) satisfies (E), this is an isomorphism of
\(\mathcal O_S\)-modules.
\footnotemark[38]

\subsection*{Remark 3.11.5.}
\footnote{Editor's note: this remark has been added.}
Suppose that \(X/S\) satisfies (E). Then
\[
  \prod_{X/S}T_{X/S}(\mathcal M)
  =\underline{\operatorname{Hom}}_{X/S}(X,T_{X/S}(\mathcal M))
\]
is a module over \(\prod_{X/S}\mathcal O_X\). Explicitly, for every
\(S'\to S\),
\begin{multline*}
  \underline{\operatorname{Hom}}_{X/S}
    (X,T_{X/S}(\mathcal M))(S')\\
  =\{\psi\in\operatorname{Hom}_X(I_{S'}(\mathcal M)\times_SX,X)
    \mid\psi\circ(\varepsilon_{\mathcal M}\times\operatorname{id}_X)
      =\operatorname{pr}_X\}
\end{multline*}
is functorially an \(\mathcal O(X\times_SS')\)-module. This follows either
from 3.6 and the properties of \(\prod_{X/S}\) (cf. 1.2), or from the
proof of 3.11.1.

We now give a geometric interpretation of the tangent bundle. Let \(U\)
be an \(S\)-functor.
\footnote{Editor's note: the original has been simplified in what
follows.}
By I~1.7.2, there are isomorphisms, functorial in
\(\mathcal M\),
\begin{align*}
  T_{X/S}(\mathcal M)(U)
  &=\operatorname{Hom}_S
    \bigl(U,\underline{\operatorname{Hom}}_S(I_S(\mathcal M),X)\bigr)\\
  &\simeq\operatorname{Hom}_S
    \bigl(I_S(\mathcal M),\underline{\operatorname{Hom}}_S(U,X)\bigr)\\
  &\simeq\operatorname{Hom}_{I_S(\mathcal M)}
    \bigl(U_{I_S(\mathcal M)},X_{I_S(\mathcal M)}\bigr).
\end{align*}
In particular, the morphism \(\mathcal M\to0\) gives the following
commutative diagram.

\SGAThreeDiagram[0.63\textwidth]
  {figures/exp2/exp2_p019_geometric_tangent_hom_square.png}
  {The geometric interpretation of the projection \(\pi_{\mathcal M}\).}
  {fig:exp2-geometric-tangent-hom}
  {Polo--Gille Exp2-14oct24.pdf, component-local p. 19; combined-reader p. 77}

The second vertical arrow is obtained by base change along
\(\varepsilon_{\mathcal M}:S\to I_S(\mathcal M)\).
\footnote{Editor's note: via
\[
  \operatorname{Hom}_{I_S(\mathcal M)}
    (U_{I_S(\mathcal M)},X_{I_S(\mathcal M)})
  \simeq\operatorname{Hom}_S(U\times_SI_S(\mathcal M),X),
\]
this is equivalent to Remark~3.1.1. Likewise, 3.12 is equivalent to the
first part of Remark~3.4.2.}

% Source page: Exp2 component-local p. 20; combined-reader p. 78.
\subsection*{Proposition 3.12.}
Let \(h_0:U\to X\) be an \(S\)-morphism. Then
\(\operatorname{Hom}_X(U,T_{X/S}(\mathcal M))\) identifies with the set of
\(I_S(\mathcal M)\)-morphisms
\[
  U_{I_S(\mathcal M)}\longrightarrow X_{I_S(\mathcal M)}
\]
whose restriction to \(U\) is \(h_0\), where \(U\) is regarded as a
subobject of \(U\times_SI_S(\mathcal M)\) through
\(\operatorname{id}_U\times_S\varepsilon_{\mathcal M}\).

Taking \(U=X\) and \(h_0=\operatorname{id}_X\) gives the following.
\footnote{Editor's note: the original has been expanded in what follows.}

\subsection*{Corollary 3.12.1.}
The set \(\Gamma(T_{X/S}(\mathcal M)/X)\) identifies with the set of
\(I_S(\mathcal M)\)-endomorphisms \(\phi\) of
\(X_{I_S(\mathcal M)}\) inducing the identity on \(X\); equivalently,
the following diagram commutes.
\footnote{Editor's note: when \(\mathcal M=\mathcal O_X\), these are
called \emph{infinitesimal endomorphisms} of \(X\); see 6.2 at the end of
this expos\'e.}

\SGAThreeDiagram[0.27\textwidth]
  {figures/exp2/exp2_p020_infinitesimal_endomorphism_square.png}
  {An infinitesimal endomorphism inducing the identity on \(X\).}
  {fig:exp2-infinitesimal-endomorphism}
  {Polo--Gille Exp2-14oct24.pdf, component-local p. 20; combined-reader p. 78}

On the other hand, by 3.11.2,
\[
  \Gamma(T_{X/S}(\mathcal M)/X)
  \simeq
  \underline{\operatorname{Lie}}
    (\underline{\operatorname{End}}_S(X)/S,\mathcal M)(S).
\]
If \(X/S\) satisfies (E), then
\(\underline{\operatorname{End}}_S(X)/S\) satisfies (E), so the latter
Lie functor is an \(\mathcal O_S\)-module by 3.6 (indeed, by 3.11.5, it
is a \(\prod_{X/S}\mathcal O_X\)-module). Applying 3.9 gives the next
proposition.
\footnote{Editor's note: the original has been expanded in what follows.}

\subsection*{Proposition 3.13.}
If \(X/S\) satisfies (E), the abelian group
\(\Gamma(T_{X/S}(\mathcal M)/X)\) identifies with the set of
\(I_S(\mathcal M)\)-endomorphisms of \(X_{I_S(\mathcal M)}\) that induce
the identity on \(X\). Consequently, every such endomorphism is an
automorphism.

\subsection*{Corollary 3.13.1.}
Let \(u:X\to Y\) be an \(S\)-isomorphism, with \(Y/S\) satisfying (E).
Every \(I_S(\mathcal M)\)-morphism
\(X_{I_S(\mathcal M)}\to Y_{I_S(\mathcal M)}\) extending \(u\) is an
isomorphism.

\subsection*{Corollary 3.13.2.}
If \(Y/S\) satisfies (E), the monomorphism
\[
  \underline{\operatorname{Isom}}_S(X,Y)
  \longrightarrow\underline{\operatorname{Hom}}_S(X,Y)
\]
induces, for every \(u\in\underline{\operatorname{Isom}}_S(X,Y)\), an
isomorphism
\[
  L^u_{\underline{\operatorname{Isom}}_S(X,Y)/S}(\mathcal M)
  \xrightarrow{\sim}
  L^u_{\underline{\operatorname{Hom}}_S(X,Y)/S}(\mathcal M).
\]

\emph{Proof.}
\footnote{Editor's note: the following proof has been added.}
For every \(S'\to S\), we must show that this map is a bijection on
\(S'\)-points. By base change (cf. 3.4), it suffices to take \(S'=S\).
The right-hand side then parametrizes the
\(I_S(\mathcal M)\)-morphisms
\(X_{I_S(\mathcal M)}\to Y_{I_S(\mathcal M)}\) extending \(u\), whereas
the left-hand side parametrizes those that are isomorphisms. The assertion
is therefore Corollary~3.13.1.

% Source page: Exp2 component-local p. 21; combined-reader p. 79.
\subsection*{Corollary 3.13.3.}
If \(X/S\) satisfies (E), the monomorphism
\(\underline{\operatorname{Aut}}_S(X)\to
\underline{\operatorname{End}}_S(X)\) induces, for every
\(u\in\underline{\operatorname{Aut}}_S(X)\), an isomorphism
\[
  L^u_{\underline{\operatorname{Aut}}_S(X)/S}(\mathcal M)
  \xrightarrow{\sim}
  L^u_{\underline{\operatorname{End}}_S(X)/S}(\mathcal M).
\]
In particular,
\[
  \underline{\operatorname{Lie}}
    (\underline{\operatorname{Aut}}_S(X)/S,\mathcal M)
  \xrightarrow{\sim}
  \underline{\operatorname{Lie}}
    (\underline{\operatorname{End}}_S(X)/S,\mathcal M)
  \xrightarrow{\sim}
  \prod_{X/S}T_{X/S}(\mathcal M),
\]
so that
\(\underline{\operatorname{Lie}}
  (\underline{\operatorname{Aut}}_S(X)/S,\mathcal M)\)
is a \(\prod_{X/S}\mathcal O_X\)-module.
\footnote{Editor's note: the preceding assertion has been added.}

\subsection*{3.14.}
Finally, suppose that \(X\) is representable.
\footnote{Editor's note: the reminder of 2.2.2 has been added, and the
sequel has been expanded.}
By 2.2.2, the \(X\)-functor \(T_{X/S}\) is represented by
\(V(\Omega^1_{X/S})\), whence bijections
\[
  \Gamma(T_{X/S}/X)
  \simeq\operatorname{Hom}_X(\Omega^1_{X/S},\mathcal O_X)
  \simeq\operatorname{Der}_{\mathcal O_S}(\mathcal O_X).
\]
This also follows from the preceding results. By 3.13,
\(\Gamma(T_{X/S}/X)\) identifies with the set of infinitesimal
endomorphisms of \(X\), that is, the \(I_S\)-endomorphisms of
\(X_{I_S}\) that induce the identity on \(X\). The schemes \(X\) and
\(X_{I_S}\) have the same underlying topological space, and their
structure sheaves are \(\mathcal O_X\) and
\[
  \mathscr D_{\mathcal O_X}=\mathcal O_X\oplus\mathcal M,
  \qquad \mathcal M=\mathcal O_X,
\]
where \(\mathcal M\) is a square-zero ideal. Let
\(\pi:\mathscr D_{\mathcal O_X}\to\mathcal O_X\) be the morphism of
\(\mathcal O_X\)-algebras vanishing on \(\mathcal M\). Giving an
infinitesimal endomorphism of \(X\) is equivalent to giving a morphism of
\(\mathcal O_S\)-algebras
\(\phi:\mathcal O_X\to\mathscr D_{\mathcal O_X}\) satisfying
\(\pi\circ\phi=\operatorname{id}_{\mathcal O_X}\), hence to giving an
\(\mathcal O_S\)-derivation of \(\mathcal O_X\).

If \(D,D'\in\operatorname{Der}_{\mathcal O_S}(\mathcal O_X)\), and
\(\phi_D\) denotes the infinitesimal endomorphism corresponding to \(D\),
then
\[
  \phi_{D+D'}=\phi_D\circ\phi_{D'}.
\]
Thus the identification
\[
  \{\text{infinitesimal endomorphisms of }X\}
  \simeq\operatorname{Der}_{\mathcal O_S}(\mathcal O_X)
\]
is an isomorphism of abelian groups. In view of 3.13 and 3.11.5, we have
therefore constructed an isomorphism of abelian groups, indeed of
\(\mathcal O(X)\)-modules,
\[
  \Gamma(T_{X/S}/X)
  \xrightarrow{\sim}\operatorname{Der}_{\mathcal O_S}(\mathcal O_X),
\]
recovering the classical interpretation of tangent vector fields as
derivations of the structure sheaf.
\footnote{Editor's note: this isomorphism is also functorial in \(S\). If
\(S'\to S\) and \(X'=X_{S'}\), then
\[
  \underline{\operatorname{Lie}}
    (\underline{\operatorname{Aut}}_S(X)/S)(S')
  \simeq\Gamma(T_{X'/S'}/X')
  \simeq\operatorname{Der}_{\mathcal O_{S'}}(\mathcal O_{X'}).
\]}
Finally,
\[
  \Gamma(T_{X/S}/X)=H^0(X,\mathfrak g_{X/S}),
\]
where \(\mathfrak g_{X/S}\) is the dual of \(\Omega^1_{X/S}\).

% Section 4 begins on Exp2 component-local p. 22; combined-reader p. 80.
% It is intentionally excluded from this component.
